\crefalias{section}{appendix}
\section{Geometric Identifiability and the Class Number}
\label{app:unitarity-class-number}
Fix an imaginary quadratic order $\mathcal O$ and let $\mathcal T_{\mathcal O}$ denote the nonempty set of complex-isomorphism classes of elliptic curves with endomorphism ring isomorphic to $\mathcal O$. The standard CM classification gives a correspondence with proper invertible ideal classes and hence $|\mathcal T_{\mathcal O}|=h(\mathcal O)$\cite{silverman_advanced_1994,waldschmidt_cm_2008}. This classification is an external mathematical input.

\begin{theorem}[Identifiability from ring data alone]
If the observation records only the fixed endomorphism-ring data, then it identifies an element of $\mathcal T_{\mathcal O}$ uniquely if and only if $h(\mathcal O)=1$.
\end{theorem}
\begin{proof}
The observation map is constant on $\mathcal T_{\mathcal O}$. A constant map is injective exactly when its domain has at most one element. The domain is nonempty, so this is equivalent to having exactly one element, which the classification identifies with $h(\mathcal O)=1$.
\end{proof}
Lean checks the abstract constant-map statement as \texttt{constant\_observation\_iff}. For the clock, Lean does not need the general correspondence: in the window $32\le|D|\le64$ it proves by Gauss reduction that all positive forms of discriminant $D$ are equivalent exactly when $D=-43$ (\texttt{clock\_unique\_iff}).

To reconstruct a physical history one must additionally specify a state space, update law, initial or final data, elapsed time, and an observation map adequate to distinguish the relevant states. A unique torus can support many trajectories. Conversely, reversible evolution is possible within each of several nonisomorphic tori. Class number one is therefore the statement that the clock's geometry is determined by its data, Postulate C; it is distinct from reversibility of the update or reconstruction of a history, which are separate requirements.
